% GENERATED by manuscript/code/build_tables.py -- do not edit by hand.
\begin{table*}[t]
\centering
\small
\caption{Pairwise comparison of baseline Brier scores on the temporal holdout using a paired sequence-cluster bootstrap.}
\label{tab:pairwise}
\begin{tabular}{llccc l}
\toprule
Model A & Model B & A $-$ B & 95\% CI & $P$(A better) & Conclusion \\
\midrule
B0 & B1 & 0.0339 & [0.0171, 0.0511] & 0.000 & resolved \\
B0 & B2 & 0.0515 & [0.0342, 0.0681] & 0.000 & resolved \\
B0 & B4 & 0.0455 & [0.0307, 0.0598] & 0.000 & resolved \\
B0 & B5 & 0.0507 & [0.0319, 0.0682] & 0.000 & resolved \\
B1 & B2 & 0.0176 & [0.0101, 0.0249] & 0.000 & resolved \\
B1 & B4 & 0.0115 & [0.0023, 0.0205] & 0.009 & resolved \\
B1 & B5 & 0.0168 & [0.0055, 0.0273] & 0.001 & resolved \\
B2 & B4 & -0.0060 & [-0.0118, 0.0002] & 0.970 & unresolved \\
B2 & B5 & -0.0008 & [-0.0078, 0.0068] & 0.613 & unresolved \\
B4 & B5 & 0.0052 & [0.0003, 0.0099] & 0.017 & resolved \\
\bottomrule
\end{tabular}
\par\vspace{2pt}\footnotesize\emph{Note.} Lower Brier score is better, so a negative A $-$ B difference favors model A. ``Resolved'' means that the 95\% interval excludes zero; ``unresolved'' means that it includes zero. Because no numerical equivalence margin was specified before analysis, unresolved comparisons are not labelled equivalent. B0 is the no-skill reference; B1 is magnitude-only logistic regression; B2 is source-only logistic regression; B4 is random forest; and B5 is XGBoost.
\end{table*}
